\documentclass{../konspekt}
%\usepackage[fontsize=5pt]{scrextend}

\title{Konspekt Termodynamika}

\begin{document}

\begin{multicols}{2}

  \section{Różniczki}

  Dla $f(x_1, \dots, x_n)$, możemy zapisać, że:
  $$
  df = \pderiv{f}{x_1} dx_1 + \dots + \pderiv{f}{x_n} dx_n
  $$

  Jeśli dla $dz(x, y) = M(x, y)dx + N(x, y)dy$;
  $$
  \pderiv{M}{y} = \pderiv{M}{x}
  $$
  to możemy powiedzieć, że $z(x, y)$ jest funkcją stanu

  Jeśli funkcja jest gładka to:
  $$
  \frac{\partial^2 z}{\partial x \partial y} = \frac{\partial^2
  z}{\partial y \partial x}
  $$
  \vspace{-2em}

  \section{Stany}

  \begin{itemize}
    \item Izotermiczny - stała temperatura \\
      $dU = 0 \quad \partial W = - \partial Q$
    \item Izobaryczny - stałe ciśnienie \\
      $dp = 0, \quad dU = \partial Q, \quad \partial Q = dH$
    \item Izochoryczny - objętość jest stała\\
      $\partial W = 0, \quad dU = \partial Q, \quad dV = 0$
    \item Adiabatyczny - ciepło jest stałe\\
      $\partial Q = 0, \quad dU = \partial W$
  \end{itemize}
  \vspace{-2em}

  \section{Pierwsza zasada termodynamiki}
  $$
  dU = \partial W + \partial Q
  $$
  Dla procesu quasi-statycznego:
  $$
  \partial W = -pdV
  $$
  \vspace{-4em}

  \section{Ciepło}
  $$
  \tcap_x = \pderiv{Q}{T}_x \quad \tcap_p = \pderiv{H}{T}_p \quad
  \tcap_V = \pderiv{U}{T}_V \quad C_x = \frac{\tcap_x}{m}
  $$
  \vspace{-2em}

  \section{Gaz Doskonały}
  $$
  pV = nRT
  $$
  $$
  \tcap_p - \tcap_V = nR \quad \tcap_V = \frac{\sigma}{2}R
  $$
  $$
  S(V, T) = C_0 + NK_B\ln(BT^{\frac{3}{2}})
  $$
  Dla $\tcap_V$ niezależnego od temperatury:
  $$
  U(T) = n\tcap_VT
  $$
  \vspace{-3em}

  \subsection{Adiabata}
  $$
  pV^\kappa = \tilde{C} = p_1V_1^\kappa = p_2V_2^\kappa \quad
  \kappa = \frac{\tcap_p}{\tcap_V} > 1
  $$
  \vspace{-2em}

  \subsection{Rozkład Maxwella}
  $$
  \rho(v) = 4 \pi \left(\frac{m_0}{2\pi k_BT}\right)^{\frac{3}{2}} v^2
  e^{-\frac{m_0v^2}{2k_BT}}
  $$
  \vspace{-2em}

  \subsection{Relacje Maxwella}
  \begin{multicols}{2}
    \begin{enumerate}
      \item $\pderiv{T}{V}_S = -\pderiv{p}{S}_V$
      \item $\pderiv{T}{p}_S = \pderiv{V}{S}_p$
      \item $\pderiv{S}{V}_T = \pderiv{p}{T}_V$
      \item $\pderiv{S}{p}_T = - \pderiv{V}{T}_p$
    \end{enumerate}
  \end{multicols}
  \vspace{-2em}

  \section{Funkcje potencjału}
  $$
  dS = \frac{\partial Q}{T} \quad S = k_B \ln T
  $$
  $$
  H(S, p) = U + pV \quad F(T, V) = U - TS \quad G(T, p) = H - TS
  $$
  $$
  \pderiv{F}{T}_V = -S = \pderiv{G}{T}_p \quad \pderiv{G}{p}_T = V
  $$
  \vspace{-2em}

  \section{Cykle}

  \begin{itemize}
    \item Carnota: $$
      \eta = \frac{\overline{W}}{Q_{AB}} = 1 - \frac{T_1}{T_2} \quad
      \overline{W} = T_1 + T_2
      $$
    \item Otta:
      $$
      Q_{DA} = n \tcap_V (T_A - T_D) > 0 \quad Q_{BC} = n \tcap_V
      (T_C - T_B) > 0
      $$
      $$
      \eta = 1 - \frac{T_B - T_C}{T_A - T_B} = 1 - r^{1 - \kappa}
      \quad r = \frac{V_C}{V_D} > 1
      $$
    \item Diesel'a:
      $$
      r = \frac{V_D}{V_A} \quad k = \frac{V_B}{V_A}
      $$
      $$
      Q_{CD} = n\tcap_V(T_D - T_C) < 0 \quad Q_{AB} = n\tcap_p (T_B - T_A) > 0
      $$
      $$
      \eta = 1 - \frac{\tcap_V(T_C - T_D)}{\tcap_p (T_B - T_A)} = 1 -
      \frac{(T_C - T_D)}{\kappa (T_B - T_A)}
      $$
  \end{itemize}
  \vspace{-2em}

  \section{Potencjał chemiczny}

  $$
  \hspace{-3em}
  \begin{aligned}
    dU = \pderiv{U}{S}_{V, n} dS + \pderiv{U}{V}_{S, n} dV +
    \pderiv{U}{n}_{S, V} dn &\rightarrow& \mu = \pderiv{U}{n}_{S, V} \\
    dF = -SdT - pdV + \mu dn &\rightarrow& \mu = \pderiv{F}{n}_{T, V} \\
    dH = TdS + Vdp + \mu dn &\rightarrow& \mu = \pderiv{H}{n}_{s, p} \\
    dG = Vdp - SdT + \mu dn & \rightarrow & \mu = \pderiv{G}{n}_{T,P}
  \end{aligned}
  $$
  \vspace{-2em}

  \section{Druga i Trzecia zasada termodynamiki}

  $$
  dS = \frac{\partial Q}{T} \quad \lim_{T \to 0} S = 0
  $$
  \vspace{-3em}

  \section{Praca Minimalna}

  $$
  W \geq \Delta U - T^{(o)} \Delta S
  $$

  \resizebox{\linewidth}{!}{%
    \centering
    \begin{tabular}{l|c|c|c|c}
      \textbf{Proces} & \textrm{const} & $\min W$ & równowaga
      & $\min$ w równowadze \\
      \hline
      izoentropowy & $S$ & $W \geq \Delta U$ & $0 \geq \Delta U$ &
      $U$ \\ \hline
      izoentropowo-choryczny & $S, V$ & $\tilde{W} \geq \Delta U$ & $0
      \geq \Delta U$ & $U$ \\ \hline
      izoentropowo-baryczny & $S, p$ & $\tilde{W} \geq \Delta H$ & $0
      \geq \Delta H$ & $H$ \\ \hline
      izotermiczny & $T$ & $W \geq \Delta F$ & $0 \geq \Delta F$ &
      $F$ \\ \hline
      izotermiczno-choryczny & $T, V$ & $\tilde{W} \geq \Delta F$ & $0
      \geq \Delta F$ & $F$ \\ \hline
      izotermiczno-baryczny & $T, p$ & $\tilde{W} \geq \Delta G$ & $0
      \geq \Delta G$ & $G$
    \end{tabular}
  }
  \vspace{-2em}

  \section{Fizyka Statystyczna}

  $$
  \frac{1}{T} = \frac{\partial S}{\partial E} \quad \mathbb{P}(Q) =
  \frac{1}{\Gamma} \quad   \beta = \frac{1}{k_BT}
  $$
  \hrule
  $$
  \mathbb{P}(Q) = \frac{\Gamma_o}{\Gamma_c}
  \quad \mathbb{P}(Q) = A \exp\left[-\beta E(Q)\right]
  $$
  \hrule
  $$
  \mathbb{P}(Q, N) = \frac{\exp \left[-\beta(E(Q_N) - \mu
  N)\right]}{\sum_N \sum_{Q_N} \mathbb{P}(Q_N, N)}
  $$

  \section{Równania Gibbsa-Durbena}

  $$
  -SdT + Vdp = \sum_i^\alpha n_i d\mu \quad \alpha + 2 - \beta = 0
  $$

  \section{Wielki potencjał termodynamiczny}

  $$
  \Omega = F - \mu n \quad d\Omega = -SdT - pdV - nd\mu
  $$

  \section{Gaz van der Waalsa}

  $$
  \left(p + \frac{n^2a}{V^2}\right)\left(V - nb\right) = nRT
  $$

\end{multicols}

\end{document}
